\section{Worked Implementation}

\begin{example}[Rebuilding post-dinner study entry as a transition design]
Consider a reader who repeatedly intends to do technical reading after dinner
but instead enters a phone-browsing state that lasts the rest of the evening.

\textbf{Step 1: Diagnose the failure.}
\begin{itemize}[nosep]
  \item Current state: post-dinner browsing with rapid novelty and low effort.
  \item Target state: twenty minutes of annotated reading with one written
        question.
  \item Dominant barrier terms: high current reward, high ambiguity about the
        first step, and a narrow post-cleanup window.
\end{itemize}

\textbf{Step 2: Identify the evidence lane.}
\begin{itemize}[nosep]
  \item Strong lane support: state-transition research justifies treating the
        dinner-to-evening boundary as a meaningful regime-change opportunity.
  \item Medium lane support: metastability and attractor-style reasoning
        support the expectation that repeated use of the same cue and same desk
        setup can make future entry more reliable.
  \item Design inference: the exact configuration below is an applied design,
        not a direct citation result.
\end{itemize}

\textbf{Step 3: Redesign the transition.}
\begin{enumerate}[nosep]
  \item \textbf{Window capture:} define the intervention point as the two
        minutes immediately after dish cleanup.
  \item \textbf{Environmental re-weighting:} phone leaves the room before the
        meal starts, so the old evening attractor is harder to reconstruct.
  \item \textbf{State preloading:} the book is opened in advance with a sticky
        note stating one question to answer.
  \item \textbf{Minimum viable entry:} read one page and write one margin note.
  \item \textbf{History shaping:} use the same chair, same lamp, and same first
        notebook action for seven consecutive sessions.
  \item \textbf{Verification:} record whether the first margin note appears
        before 8:15 PM.
\end{enumerate}

\textbf{Why this is better than generic advice.}
The design does not ask the reader to feel more serious. It changes the local
geometry of the transition. The dinner boundary becomes a designed decision
window; the target state is partially built before entry; the old state is made
less recoverable; and repeated history is used to increase the odds that the
next evening begins from a more favourable local configuration.
\end{example}
